% UC09 -- Calcular Correlação Sono × Energia

\begin{longtable}{| p{.40\textwidth} | p{.50\textwidth} |}
\hline

\textbf{Nome do Caso de Uso} &
UC09 -- Calcular Correlação Sono $\times$ Energia \\
\hline

\textbf{Ator Principal} &
Sistema \\
\hline

\textbf{Atores Secundários} &
--- \\
\hline

\textbf{Resumo} &
Calcula automaticamente a correlação entre os registros recentes de sono e os níveis de energia informados pelo usuário. \\
\hline

\textbf{Pré-condições} &
Existem registros suficientes para o cálculo da correlação. \\
\hline

\textbf{Pós-condições} &
O indicador de correlação é atualizado. \\
\hline

\multicolumn{2}{|c|}{\textbf{Cenário Principal}}\\
\hline

\textbf{Ações do Ator} &
\textbf{Resposta do Sistema} \\
\hline

&
1. O sistema identifica a criação ou atualização de um registro de humor ou sono. \\
\cline{2-2}

&
2. O sistema recupera os registros correspondentes aos últimos sete dias. \\
\cline{2-2}

&
3. O sistema calcula a correlação utilizando o coeficiente de Pearson. \\
\cline{2-2}

&
4. O sistema atualiza o indicador apresentado ao usuário. \\
\hline

\multirow{2}{.40\textwidth}{\textbf{Restrições e validações}}
&
O cálculo considera apenas registros dos últimos sete dias. \\
\cline{2-2}

&
É necessário existir quantidade mínima de registros para o cálculo estatístico. \\
\hline

\multicolumn{2}{|c|}{\textbf{Cenário Alternativo A1 -- Dados insuficientes}}\\
\hline

&
O sistema mantém o indicador indisponível até que existam registros suficientes para o cálculo. \\
\hline

\end{longtable}
